{\footnotesize
    \setlength{\tabcolsep}{4pt}
    \renewcommand{\arraystretch}{1.10}
    \setlength{\LTleft}{0pt}
    \setlength{\LTright}{0pt}
    \setlength{\LTcapwidth}{\linewidth}
    \newlength{\PenelitianTerdahuluWidth}
    \setlength{\PenelitianTerdahuluWidth}{\dimexpr\linewidth-6\tabcolsep-3pt\relax}
    \refstepcounter{table}\label{tab:penelitian-terdahulu}
    \edef\PenelitianTerdahuluNomor{\thetable}
    \addcontentsline{lot}{table}{\protect\numberline{\thetable}Penelitian Terdahulu yang Relevan}
    \begin{longtable}{@{}>{\raggedright\arraybackslash}p{0.14\PenelitianTerdahuluWidth}>{\raggedright\arraybackslash}p{0.30\PenelitianTerdahuluWidth}>{\raggedright\arraybackslash}p{0.18\PenelitianTerdahuluWidth}>{\raggedright\arraybackslash}p{0.38\PenelitianTerdahuluWidth}@{}}
        \multicolumn{4}{@{}c@{}}{\textbf{Tabel \PenelitianTerdahuluNomor}}\\
        \multicolumn{4}{@{}c@{}}{\textbf{Penelitian Terdahulu yang Relevan}}\\[6pt]
        \toprule
        \textbf{Nama (Tahun)}                     & \textbf{Judul Penelitian}                                                                                                                                                            & \textbf{Metode Analisis}                                                                                                                       & \textbf{Hasil Penelitian}                                                                                                                                                                                                                                           \\
        \midrule
        \endfirsthead
        \multicolumn{4}{@{}l@{}}{\textbf{Lanjutan Tabel \PenelitianTerdahuluNomor}}\\
        \toprule
        \textbf{Nama (Tahun)}                     & \textbf{Judul Penelitian}                                                                                                                                                            & \textbf{Metode Analisis}                                                                                                                       & \textbf{Hasil Penelitian}                                                                                                                                                                                                                                           \\
        \midrule
        \endhead
        \midrule
        \multicolumn{4}{r}{\footnotesize Dilanjutkan pada halaman berikutnya}\\
        \endfoot
        \bottomrule
        \endlastfoot
        \textcite{regina2025kuliner}              & Strategi Pemasaran Digital untuk Meningkatkan Penjualan UMKM Sektor Kuliner                                                                                                          & Deskriptif kualitatif                                                                                                               & Strategi digital melalui media sosial, katalog daring, dan pemesanan online meningkatkan visibilitas usaha, interaksi dengan konsumen, dan peluang kenaikan penjualan UMKM kuliner.                                                      \\
        \textcite{sophia2025rendanggadih}         & Analisis Strategi Pemasaran Rendang Gadih Melalui Platform Shopee untuk Meningkatkan Penjualan                                                                                       & Deskriptif kualitatif                                                                                                               & Peningkatan penjualan dicapai melalui penerapan STP, bauran pemasaran 4P, serta pemanfaatan fitur Shopee seperti CRM, Live, Ads, dan Campaign.                                                   \\
        \textcite{hendrizon2024madayacoffee}      & Analisis Strategi Pemasaran pada Madaya Coffee Parung Bogor                                                                                                                          & Matriks IFE, EFE, SWOT, dan QSPM                                                                                                    & Madaya Coffee perlu meningkatkan aktivitas mengikuti pameran untuk memperkenalkan produk kepada khalayak lebih luas sebagai bagian dari upaya peningkatan penjualan.                              \\
        \textcite{ramadhanti2023rucika}           & Analisis SWOT dalam Menentukan Strategi Pemasaran Rucika pada Distributor CV Sinar Mas Agung Semarang                                                                                & Kualitatif dengan studi kasus; wawancara, observasi, dokumentasi, SWOT, dan bauran pemasaran 7P                         & Strategi pemasaran CV Sinar Mas Agung mengarah pada \textit{growth oriented strategy}; aspek promosi dan bukti kepuasan konsumen masih perlu dimaksimalkan.                                      \\
        \textcite{haliza2023seqapla}              & Analisis Strategi Pemasaran dalam Upaya Peningkatan Penjualan pada Seqapla Coffee                                                                                                    & Analisis deskriptif kondisi internal-eksternal dan perumusan strategi pemasaran                                                     & Seqapla Coffee membutuhkan penguatan strategi pemasaran yang lebih terarah untuk menghadapi persaingan \textit{coffee shop} dan mendorong pemulihan penjualan.                                  \\
        \textcite{ruhiyat2023minatbeli}           & Strategi Pemasaran Online untuk Meningkatkan Minat Beli Konsumen Produk Makanan dan Minuman UMKM di Kota Bogor                                                                       & Regresi linear berganda, analisis deskriptif, SWOT, IFE-EFE, IE, dan QSPM                                                           & Unsur 4P berpengaruh positif terhadap minat beli; posisi UMKM berada pada strategi \textit{grow and build} dengan rekomendasi inovasi produk sesuai kebutuhan dan daya beli konsumen.             \\
        \textcite{khairani2022buahsayur}          & Strategi Meningkatkan Penjualan Buah dan Sayur pada Platform Digital di Kota Bogor                                                                                                   & Analisis deskriptif, regresi data panel, dan SWOT                                                                                   & Harga, ulasan, peringkat, jenis toko, dan promosi memengaruhi penjualan; strategi yang direkomendasikan adalah kerja sama langsung dengan produsen sebagai pemasok utama.                        \\
        \textcite{mamuriyah2021gadihminang}       & Strategi Pemasaran dalam Meningkatkan Hasil Penjualan Rumah Makan Padang Gadih Minang                                                                                                & Kegiatan PKM dengan promosi digital melalui brosur, Instagram, dan YouTube                                                          & Promosi melalui brosur, Instagram, dan YouTube diarahkan untuk memperkenalkan Rumah Makan Padang Gadih Minang ke khalayak yang lebih luas dan diharapkan mendorong kenaikan omzet penjualan.    \\
        \textcite{roslam2021markisaaurora}        & Strategi Pemasaran UMKM Markisa Aurora di Kota Makassar pada Masa Pandemi COVID-19                                                                                                   & Analisis deskriptif strategi pemasaran UMKM pada masa pandemi                                                                       & UMKM Markisa Aurora dapat mempertahankan dan memperluas pasar melalui penguatan promosi digital, penyesuaian strategi pemasaran, dan pemanfaatan kanal penjualan daring selama pandemi.          \\
        \textcite{suharjo2020bogorsnacks}         & Strategi Pemasaran Digital pada UMKM Makanan Ringan di Kota Bogor                                                                                                                    & Analisis deskriptif, perilaku pembelian daring konsumen, \textit{Importance Performance Analysis}, bauran pemasaran digital, dan 4C & Strategi prioritas menekankan penguatan konten digital, pemanfaatan media sosial, dan penyesuaian elemen bauran pemasaran digital untuk memperluas jangkauan pasar UMKM makanan ringan.        \\
        \textcite{gurel2017swot}                  & SWOT Analysis: A Theoretical Review                                                                                                                                                  & Tinjauan teoretis terhadap konsep dan penggunaan SWOT                                                                               & Kajian ini menegaskan bahwa SWOT merupakan alat analisis strategis yang membantu organisasi mengidentifikasi kekuatan, kelemahan, peluang, dan ancaman sebagai dasar perumusan strategi.         \\
        \textcite{helms2010swotreview}            & Exploring SWOT Analysis -- Where Are We Now?: A Review of Academic Research from the Last Decade                                                                                     & Kajian pustaka terhadap penelitian SWOT dalam satu dekade                                                                           & Hasil telaah menunjukkan bahwa SWOT tetap luas digunakan dalam penelitian strategi, tetapi memerlukan penerapan yang lebih sistematis dan dukungan kerangka analitis yang lebih kuat.             \\
    \end{longtable}
    \tablesource[\linewidth]{Diolah peneliti dari berbagai penelitian terdahulu (2026).}
    \addtocounter{table}{-1}
}
